\item \textbf{ACTIVIDAD} Para esta actividad, necesitará algunos recipientes, algunos popotes y un poco de agua. 
(a) Llene varios recipientes a diferentes niveles con agua y mida la profundidad $h$ del agua en cada taza cuidadosamente. Ahora sumerja un popote verticalmente en cada taza con su parte inferior apoyándose en la parte inferior de la taza. Coloque su dedo sobre el extremo superior del popote para sellar y levantarlo verticalmente fuera de la taza. Mida la longitud $h'$ de la columna de agua atrapada en el popote. Es posible que tenga que tomar una foto del popote con el teléfono inteligente y una regla y analizar cuidadosamente una ampliación de la foto para hacer esta medida. ¿La longitud de la columna de agua atrapada en el popote coincide con la profundidad del agua en la taza? ¿Debería coincidir? 
(b) Para que la columna de agua se suspenda en el popote, a presión sobre la columna (dentro de la paja) debe ser menor que la presión atmosférica. Para que esto suceda, la columna de agua debe moverse hacia abajo un poco cuando el popote se levanta del agua, por lo que el aire encima de la columna se expande en el volumen y la presión disminuye. Por lo tanto, las dos longitudes medidas no deben ser las mismas. Pero, ¿es detectable la diferencia? Demuestre que la longitud $h'$ del agua suspendida en el popote de longitud $\ell$ debe estar relacionada con la profundidad $h$ del agua en la taza por
$$h' = \frac{1}{2}\left(\ell + \frac{P_0}{\rho g}\right) - \frac{1}{2}\sqrt{\ell^2 + 2\left(\frac{P_0}{\rho g}\right)(\ell - 2h) + \left(\frac{P_0}{\rho g}\right)^2}$$
donde $P_0$ es la presión atmosférica y $\rho$ es la densidad del agua. 
(c) En el caso de un popote de 30 cm, utilice la ecuación para determinar la diferencia $h - h'$ para diferentes niveles de $h$ de cero a 30 cm en incrementos de 1 cm. 
(d) ¿Para qué valores de $h$ es el porcentaje de goteo en la columna de agua más grande cuando el popote se saca de la taza? 
(e) ¿Para qué valor de $h$ es el valor de $h - h'$ un máximo?
